% ============================================================================
% CHAPTER 5: CONCLUSION AND FUTURE WORK
% ============================================================================

\chapter{بحث و نتیجه‌گیری}

\section{مقدمه}
در این پژوهش، یک الگوریتم نوین برای رمزنگاری تصاویر رنگی با استفاده از ترکیب سیستم‌های آشوبی سه‌بعدی و مدل‌های آشوبی نمایی (\lr{ECM}) ارائه گردید. هدف اصلی، افزایش سطح امنیت و پیچیدگی فرآیند رمزنگاری در کنار حفظ کارایی محاسباتی بود.

\section{بحث و تفسیر نتایج}
نتایج به‌دست‌آمده نشان می‌دهد که ترکیب ۹ نقشه آشوبی نمایی به‌صورت دینامیک برای هر پیکسل، کارایی امنیتی بالایی ایجاد کرده است. به‌کارگیری سیستم \lr{Chen} برای جایگشت، همبستگی مکانی تصویر را به حداقل می‌رساند. تحلیل هیستوگرام و آنتروپی نشان می‌دهد که توزیع پیکسل‌ها به نحو قابل توجهی یکنواخت شده است و تصادفی بودن داده‌ها تضمین شده است.

مقادیر \lr{NPCR} و \lr{UACI} بیانگر قدرت اثر \lr{Avalanche} و مقاومت بالا در برابر حملات دیفرانسیلی است. همچنین، زمان اجرای کمتر از ۳ ثانیه برای تصاویر $512 \times 512$ بیانگر کارایی عملی و قابلیت استفاده در سیستم‌های بلادرنگ است.

\section{مقایسه با سایر الگوریتم‌ها}
برای تحلیل جامع‌تر، نتایج الگوریتم پیشنهادی با کارهای اخیر مقایسه شده است. الگوریتم پیشنهادی موفق شده با تنها یک دور، کارایی و امنیتی مشابه یا بهتر از روش‌های چنددوری ارائه دهد. استفاده از ۹ نقشه نمایی به‌صورت انتخاب دینامیک، مقاومت در برابر تحلیل آماری را به شکل قابل توجهی افزایش داده است.

\section{نتیجه‌گیری نهایی}
با توجه به نتایج عددی و بصری، تمامی فرضیه‌های تحقیق تأیید شده‌اند:
\begin{enumerate}
    \item افزایش امنیت اطلاعات تصویری.
    \item کاهش همبستگی پیکسل‌های مجاور به کمتر از ۰.۰۵٪.
    \item دستیابی به آنتروپی نزدیک به ۸ (۷.۹۷۸۹).
    \item کارایی محاسباتی بالا (زمان اجرا کمتر از ۳ ثانیه).
\end{enumerate}
ترکیب ۹ نقشه نمایی دینامیک و سیستم \lr{Chen} برای جایگشت، نوآوری کلیدی این پژوهش محسوب می‌شود.

\section{پیشنهادات برای پژوهش‌های آتی}
برای ادامه این پژوهش، موارد زیر پیشنهاد می‌شود:
\begin{enumerate}
    \item استخراج کلید با \lr{SHA-512} یا \lr{SHA-3} از تصویر اصلی برای افزایش امنیت.
    \item اعمال چند دور رمزنگاری با حالت زنجیره‌ای (\lr{CBC-like}) برای افزایش اثر دیفرانسیلی.
    \item گسترش الگوریتم به رمزنگاری ویدئو و تصاویر پزشکی با حجم بالا.
    \item پیاده‌سازی روی \lr{FPGA} و مقایسه مصرف توان و کارایی واقعی.
    \item ترکیب الگوریتم با روش‌های یادگیری عمیق برای رمزنگاری تطبیقی.
\end{enumerate}
\
